% !PRINTED 340
\weberpage{340}

Setzt man dann in (5)
\[
x = \alpha' \quad \text{und} \quad x = \beta'
\]
und beachtet, dass
\[
\beta - \alpha' = \frac{(\beta - \alpha)f(\beta)}{f(\beta) - f(\alpha)}, \quad \beta - \beta' = \frac{f(\beta)}{f'(\beta)},
\]
so ergiebt sich
\[
f(\alpha') < 0, \quad f(\beta') > 0.
\]

Es ist also die gesuchte Wurzel zwischen $\alpha'$ und $\beta'$ enthalten,
und das Intervall $(\alpha', \beta')$ ist kleiner als $(\alpha, \beta)$.

Setzen wir in diesen Ausdr\"ucken $\alpha', \beta'$ an Stelle von $\alpha, \beta$,
so erhalten wir ein neues, noch engeres Intervall u.\,s.\,f.

Was die Convergenz dieses Verfahrens betrifft, so k\"onnen
wir uns dar\"uber folgendermaassen vergewissern. Wir setzen
nach (6):
\weq{7}{
\frac{\beta' - \alpha'}{\beta - \alpha} = \Phi(\alpha, \beta) = \frac{f(\beta)\{f(\alpha) - f(\beta) + (\beta - \alpha)f'(\beta)\}}{(\beta - \alpha)f'(\beta)\{f(\beta) - f(\alpha)\}}.
}
Z\"ahler und Nenner sind hier durch $(\beta - \alpha)^2$ theilbar, und
wenn wir den gemeinsamen Factor $(\beta - \alpha)^2$ wegheben und $\xi, \eta$
f\"ur $\alpha, \beta$ setzen, so erhalten wir einen Ausdruck von der Form
\[
\Phi(\xi, \eta) = \frac{\psi_1(\xi, \eta)}{\psi_2(\xi, \eta)},
\]
worin $\psi_1$ und $\psi_2$ ganze rationale Functionen der beiden Variablen
$\xi, \eta$ sind.

Ist $x$ die im Intervall $(\alpha, \beta)$ gelegene Wurzel, so werden die
beiden Functionen $\psi_1$ und $\psi_2$ nach unseren Voraussetzungen
\"uber $f(x)$ nicht gleich Null, wenn
\weq{8}{
\alpha \leqq \xi < x, \quad x < \eta \leqq \beta.
}
Die Function $\Phi(\xi, \eta)$ ist f\"ur $\xi = \alpha, \eta = \beta$, aber auch f\"ur
jedes andere Werthpaar in dem Intervall $(\alpha, \beta)$, wof\"ur $f(\xi)$
negativ, $f(\eta)$ positiv ist, kleiner als 1, da jedes solche Intervall
$(\xi, \eta)$ an Stelle von $(\alpha, \beta)$ genommen werden k\"onnte.

Da nun $\Phi(\xi, \eta)$ eine stetige Function der beiden Ver\"anderlichen
 $\xi, \eta$ ist, so lange diese in dem Bereich (8) bleiben, so
muss in diesem Bereich die Function $\Phi(\xi, \eta)$ einen Maximumwerth
 haben, und dieser muss kleiner als 1 sein, weil $\Phi(\xi, \eta)$
auch noch in den Grenzf\"allen $\xi = x, \eta = x$ kleiner als 1 bleibt.

Es l\"asst sich also ein positiver echter Bruch $\Theta$ angeben,
so dass
\[
\Phi(\xi, \eta) < \Theta
\]
%CONT>>
